\frontslide{Hachage linéaire}


%------------------------------------------------------------
\begin{frame}{Le hachage linéaire}

Hachage linéaire: un hachage qui s'adapte dynamiquement aux insertions
et suppressions, avec une \alert{évolution minimale} du répertoire
et des fragments.

\vfill

\begin{tabular}{l|c}
titre & $h$(titre)   \\
\hline
Vertigo     & 14  \\
Brazil      & 43  \\
Twin Peaks  & 25  \\
Underground & 24  \\
Easy Rider  & 8 \\
Psychose    &  33 \\
Greystoke   &  17 \\
Shining     &  16 \\
Citizen Kane     &  48 \\\hline
\end{tabular}

\vfill
Une fonction $h(c)$ immuable 
s'applique à la clé et renvoie un entier.
\end{frame}
%------------------------------------------------------------
\begin{frame}{La structure}

Comme les autres hachages, plus un paramètre\,:  \alert{l'indice de partitionnement}.

\vfill

  \figSlideWithSize{hachageLineaire1}{8cm}

\vfill
\begin{blocgauche}
On considère la famille de fonctions $h_i (c)  = h(c)\, mod\, 2^i$. Ici
on appplique $h_1$ (modulo 2). 
\end{blocgauche}
\end{frame}

\begin{frame}{Insertions}

On insère Psychose, puis Greystoke qui vont tous deux dans $f_1$.
C'est $f_1$ qui déborde, \alert{mais c'est $f_0$ qui éclate en ($f_0, f_2$)}.
\vfill

  \figSlideWithSize{hachageLineaire2}{9cm}


\vfill
\begin{blocgauche}
Maintenant on a deux fonctions: $h_1$ s'applique pour $f_1$, $h_2$
pour les autres. Et $p$ vaut 1.
\end{blocgauche}
\end{frame}



\begin{frame}{On insère Shining puis Citizen Kane}

Tous les deux vont dans $f_0$ qui déborde.  \alert{Mais maintenant c'est $f_1$
qui éclate}. 
\vfill

  \figSlideWithSize{hachageLineaire3}{10cm}


\vfill

\begin{blocgauche}
On a \alert{découplé} les éclatements et les débordements. On sait que 
les seconds seront rectifiés par les premiers, \alert{à terme}. 

\smallskip

 $p$ est revenu à 0; la fonction est $h_2$: c'est reparti pour un tour.
\end{blocgauche}

\end{frame}



\begin{frame}{Résumé du hachage linéaire}

À chaque moment, on a deux fonctions $(h_{n}, h_{n+1})$ et un indice
de partitionnement $p$.

\vfill

\begin{redbullet}
  \item La fonction $h_{n+1}$ s'applique à tous les fragments \alert{avant} $p$
       car ils ont déjà éclaté.
     \item La fonction $h_{n+1}$ à tous les autres fragments qui n'ont pas
       encore éclaté.
\end{redbullet}

\vfill

Quand le dernier fragment éclate ($p=n$) la fonction $h_n$ ne sert plus et la
paire de fonction de hachage devient $(h_{n+1}, h_{n+2})$
  

\vfill

Brillant!
\end{frame}